\noindent
\begin{minipage}[t]{0.23\textwidth}
  % Cột lề trái: chứa ghi chú lề được dóng hàng chuẩn xác theo tọa độ gốc
  \vspace*{19.8cm}
  \begin{flushleft}
    \fontsize{8.5pt}{11pt}\selectfont
    Xem Trang tra cứu 2 hoặc Phụ lục D.
  \end{flushleft}
\end{minipage}%
\hfill
\begin{minipage}[t]{0.73\textwidth}
  % Cột chính: nội dung bài học, ví dụ, công thức
  \noindent
  \textbf{\textsf{\color{stewartcyan}VÍ DỤ 1}}\enspace Tìm $F'(x)$ nếu $F(x) = \sqrt{x^2 + 1}$.

  \vspace{0.2em}
  \noindent
  \textbf{\textsf{\color{stewartcyan}LỜI GIẢI 1}}\enspace (sử dụng Công thức 1): Ở phần đầu của mục này, chúng ta đã biểu diễn $F$ dưới dạng $F(x) = (f \circ g)(x) = f(g(x))$ với $f(u) = \sqrt{u}$ và $g(x) = x^2 + 1$. Vì
  \[
  f'(u) = \frac{1}{2}u^{-1/2} = \frac{1}{2\sqrt{u}} \qquad \text{và} \qquad g'(x) = 2x
  \]
  ta có
  \begin{align*}
  F'(x) &= f'(g(x)) g'(x)\\[3pt]
  &= \frac{1}{2\sqrt{x^2 + 1}} \cdot 2x = \frac{x}{\sqrt{x^2 + 1}}
  \end{align*}

  \noindent
  \textbf{\textsf{\color{stewartcyan}LỜI GIẢI 2}}\enspace (sử dụng Công thức 2): Nếu ta đặt $u = x^2 + 1$ và $y = \sqrt{u}$, khi đó
  \[
  F'(x) = \frac{dy}{du}\frac{du}{dx} = \frac{1}{2\sqrt{u}}(2x) = \frac{1}{2\sqrt{x^2 + 1}}(2x) = \frac{x}{\sqrt{x^2 + 1}} \tag*{\color{stewartcyan}\blacksquare}
  \]

  \vspace{0.3em}
  Khi sử dụng Công thức 2, chúng ta cần ghi nhớ rằng $dy/dx$ chỉ đạo hàm của $y$ khi $y$ được coi là một hàm theo biến $x$ (đạo hàm của $y$ theo biến $x$), trong khi $dy/du$ chỉ đạo hàm của $y$ khi được coi là một hàm theo biến $u$ (đạo hàm của $y$ theo biến $u$). Chẳng hạn, trong Ví dụ 1, $y$ có thể được coi là một hàm theo biến $x$ ($y = \sqrt{x^2 + 1}$) và cũng là một hàm theo biến $u$ ($y = \sqrt{u}$). Chú ý rằng
  \[
  \frac{dy}{dx} = F'(x) = \frac{x}{\sqrt{x^2 + 1}} \qquad \text{trong khi} \qquad \frac{dy}{du} = f'(u) = \frac{1}{2\sqrt{u}}
  \]

  \vspace{0.4em}
  \noindent
  \textbf{\textsf{\color{stewartgreen}CHÚ Ý}}\enspace Khi sử dụng Quy tắc chuỗi, chúng ta làm việc từ ngoài vào trong. Công thức 1 nói rằng chúng ta \textit{lấy đạo hàm của hàm ngoài $f$ [tính tại hàm trong $g(x)$] rồi nhân với đạo hàm của hàm trong}.
  \[
  \frac{d}{dx} \cunder{f}{hàm\\ngoài} \;\cunder{(g(x))}{tính tại\\hàm trong} \;=\; \cunder{f'}{đạo hàm\\của hàm ngoài} \;\cunder{(g(x))}{tính tại\\hàm trong} \;\cdot\; \cunder{g'(x)}{đạo hàm\\của hàm trong}
  \]

  \vspace{0.4em}
  \noindent
  \textbf{\textsf{\color{stewartcyan}VÍ DỤ 2}}\enspace Tính đạo hàm của (a) $y = \sin(x^2)$ và (b) $y = \sin^2 x$.

  \vspace{0.2em}
  \noindent
  \textbf{\textsf{\color{stewartcyan}LỜI GIẢI}}

  (a) Nếu $y = \sin(x^2)$, thì hàm ngoài là hàm sin và hàm trong là hàm bình phương, do đó Quy tắc chuỗi cho ta
  \begin{align*}
  \frac{dy}{dx} &= \frac{d}{dx} \cunder{\sin}{hàm\\ngoài} \;\cunder{(x^2)}{tính tại\\hàm trong} \;=\; \cunder{\cos}{đạo hàm\\của hàm ngoài} \;\cunder{(x^2)}{tính tại\\hàm trong} \;\cdot\; \cunder{2x}{đạo hàm\\của hàm trong}\\[6pt]
  &= 2x \cos(x^2)
  \end{align*}

  (b) Lưu ý rằng $\sin^2 x = (\sin x)^2$. Ở đây hàm ngoài là hàm bình phương và hàm trong là hàm sin. Vì vậy
  \[
  \frac{dy}{dx} = \frac{d}{dx} \cunder{(\sin x)}{hàm\\trong}^2 \;=\; \cunder{2}{đạo hàm\\của hàm ngoài} \;\cdot\; \cunder{(\sin x)}{tính tại\\hàm trong} \;\cdot\; \cunder{\cos x}{đạo hàm\\của hàm trong}
  \]

  \vspace{0.2em}
  Kết quả có thể để nguyên là $2\sin x\cos x$ hoặc viết thành $\sin 2x$ (theo một đồng nhất thức lượng giác quen thuộc là công thức nhân đôi). \hfill {\color{stewartcyan}\blacksquare}
\end{minipage}
